% ---- VII. Discussion. Source: JAMA Discussion paragraphs 1 and 2; limitations are new, one
% per JAMA reviewer objection. The Regional Medical Center aside is dropped.
\section{Discussion}
\label{sec:discussion}

The MDP policy improves expected outcomes over nearest-hospital routing in a decentralized system and matches the centralized protocol of a small state while spreading load across hospitals. By accounting for hospital capability, elapsed time and the cost of a later transfer, it delivers more effective routing where it matters for equity: in areas distant from EVT-capable centers, with no systematic loss in well-served areas, and the stratified analysis shows the gain over the bypass heuristics as well as over the default. The decision tree carries this benefit into a form a crew can apply, and a decision costs milliseconds.

Two practical considerations merit attention. The MDP increases mean initial transport time to 16.7 minutes from 9.4 under nearest-hospital routing in the Bay Area, but it reduces downstream transfers and their delays; for LVO patients, bypassing a closer but less capable hospital means arriving at an EVT-capable center within the therapeutic window. At the same time the MDP travels less than the nearest-EVT-center rule (21.4 minutes), because it declines the detour for patients whose outcome it cannot change. The second consideration is hospital load, discussed below.

The in-hospital intervals matter as much as the road network. With door-in-door-out at 60 minutes the gain over nearest-hospital routing is a third of its value at the registry median, and at the registry upper quartile it is larger still: the value of intelligent routing grows with the inefficiency of transfer. Agencies with fast transfer pathways should expect smaller gains, and an agency that measures its own intervals can substitute them directly.

\subsection{Limitations}
\emph{Capacity and queueing.} The model assumes every hospital can accept every patient routed to it. We do not model bed or angiography-suite availability, diversion, or interaction between concurrent patients. The load-distribution result (the MDP uses 51 hospitals where the nearest-EVT rule uses 10, and places 35\% rather than 52\% of patients at the two busiest centers) is a qualitative indication that the policy does not concentrate demand, not a capacity analysis. A capacity-aware extension would add hospital state to the MDP and is the main direction for future work.

\emph{Validation against observed outcomes.} Results are simulated. The outcome model is taken from published onset-to-treatment curves~\cite{holodinsky2018}, travel times from a routing engine, and in-hospital intervals from national registries; the composition has not been validated against a cohort with observed routing and outcomes. The sensitivity analyses bound the effect of the main inputs but do not substitute for validation, which requires EMS and hospital record linkage.

\emph{In-hospital intervals.} One value per interval for all hospitals of a tier. Real door-to-needle, door-to-puncture and door-in-door-out times vary widely between hospitals and over time; the DIDO sensitivity shows the direction and size of the effect, and hospital-specific intervals can be substituted where an agency has them.

\emph{Patient population.} The cohort is LAMS~$\geq$~4 with onset within 270 minutes. Prehospital scales identify LVO imperfectly, and EVT benefit extends to 24 hours for selected patients. The outcome curves used here are defined to 270 minutes; extending the model to the late window needs late-window outcome functions and is left for future work. Less severe strokes also face routing decisions that the present scope does not address.

\emph{Hemorrhage.} The outcome for intracerebral hemorrhage is time-invariant in the reward. Evidence that timely bundled care improves ICH outcomes~\cite{interact3} suggests a time-dependent function would be appropriate when one is available; its absence biases the model against any destination choice motivated by ICH, which is conservative for the MDP's gains.

\emph{Capability definitions.} Hospitals are classified by verified certification. Some hospitals perform thrombectomy without certification; the on-site sensitivity analysis treats them as EVT-capable and does not change the conclusions. Designations change over time; the method takes the table as input.

\emph{Generalizability and the baseline.} Both regions are dense relative to much of the United States, and routing protocols vary by jurisdiction. Nearest-hospital routing is a lower bound on current practice rather than a description of it everywhere; the heuristic comparators correspond to the protocols actually in use, which is why both are reported. Air transport is not modeled.

\emph{Travel times.} Deterministic routing-engine estimates stand in for traffic-dependent times; the noise and bias analyses cover the plausible range of error but not time-of-day structure.

\subsection{Deployment path and other applications}
The deployable artifact is the decision tree, re-derived whenever the hospital table changes, with the MDP as the reference it is checked against. The natural next step is a prospective evaluation with an EMS agency in which the tree's recommendation is recorded alongside the crew's destination and the patient's eventual outcome. The formulation is not specific to stroke: any time-critical dispatch decision with tiered receiving facilities, an unobserved patient class revealed on arrival, and a possible secondary transfer (trauma, ST-elevation myocardial infarction, aortic dissection) has the same structure and would need only its own outcome curves and facility table.
